\subsection{\texorpdfstring{\(R^{m}\) 到 \(T^{n}\) 中的连续同态}{R 的 m 次幂到 T 的 n 次幂中的连续同态}}

命题 3. \(\mathbf{R}^m\) 到 \(\mathbf{T}^n\) 中的每个连续同态都形如 \(x \to \varphi(u(x))\) ，其中 \(\varphi\) 是 \(\mathbf{R}^n\) 到 \(\mathbf{T}^n\) （与 \(\mathbf{R}^n / \mathbf{Z}^n\) 等同）上的典范同态，而 \(u\) 是 \(\mathbf{R}^m\) 到 \(\mathbf{R}^n\) 中的线性映射。

设 \(f\) 是 \(\mathbf{R}^m\) 到 \(\mathbf{T}^n\) 中的一个连续同态。我们将证明，存在线性映射 \(u\) （ \(\mathbf{R}^m\) 到 \(\mathbf{R}^n\) 中的），使得同态 \(x \to f(x)\) 与 \(x \to \varphi(u(x))\) 在 \(o\) 于 \(\mathbf{R}^m\) 中的一个邻域的所有点处一致；于是命题由第 2 小节命题 2 的推论 1 得出。现在设 \(V\) 是 \(o\) 于 \(\mathbf{R}^n\) 中的一个这样的邻域： \(\varphi\) 在 \(V\) 上的限制是 \(\mathbf{R}^n\) 与 \(T^n\) 之间的一个局部同构，并设 \(\psi\) 是 \(\varphi\) 的逆映射，定义在 \(\varphi(V)\) 上。由于 \(f\) 连续， \(V' = \overline{f}^{-1}(\varphi(V))\) 是 \(o\) 于 \(\mathbf{R}^m\) 中的一个邻域；映射

\[
\boldsymbol {x} \rightarrow \psi (f (\boldsymbol {x})),
\]

限制在 \(V'\) 上，是 \(V'\) 到 \(R^n\) 中的连续映射，使得 \(\psi(f(x+y)) = \psi(f(x)) + \psi(f(y))\) 对 \(R^m\) 中每对满足 \(x \in V'\) 、 \(y \in V'\) 与 \(x + y \in V'\) 的点 x、y 成立；因而（第 2 小节命题 2 的推论 1）这个映射与一个完全确定的 \(R^m\) 到 \(R^n\) 中的连续同态 u 在 \(\mathbf{R}^m\) 中 o 的一个邻域 W 的所有点处一致。由命题 1（第 1 小节）， \(u\) 是 \(\mathbf{R}^m\) 到 \(\mathbf{R}^n\) 中的线性映射；于是对一切 \(x \in W\) 有 \(f(x) = \varphi(u(x))\) ，证明完毕。

注. 同样的论证更一般地表明：若 \(\varphi\) 是 \(R^{n}\) 到群 G 中的严格态射，它在 o 的适当邻域上的限制是 \(R^{n}\) 与 G 之间的局部同构，则 \(R^{n}\) 到 G 中的每个连续同态都形如 \(x \to \varphi(u(x))\) ，其中 u 是 \(R^{m}\) 到 \(R^{n}\) 中的线性映射。

在 \(m = n = 1\) 的情形，命题 3 给出：

命题 4. 若 \(\varphi\) 是 \(\mathbf{R}\) 到 \(\mathbf{T}\) 上的典范同态，则 \(\mathbf{R}\) 到 \(\mathbf{T}\) 中的每个连续同态都形如 \(x \to \varphi(ax)\) ，其中 \(a \in \mathbf{R}\) ；且它是 \(\mathbf{R}\) 到 \(\mathbf{T}\) 上的严格态射（当 \(a \neq 0\) 时）。

